\documentclass[aspectratio=169,11pt]{beamer}
% Shared look for every deck: Metropolis with a progress bar under each frame title.
% Build with Tectonic: `make` in this directory.

\usetheme[progressbar=frametitle,numbering=fraction,sectionpage=progressbar,block=fill]{metropolis}

% Metropolis looks Fira Sans up by system name; load it from the TeX bundle instead,
% Light for body text as Metropolis intends, with lining figures for data.
\usepackage[sfdefault,light,lf]{FiraSans}
\usepackage{booktabs,tabularx}
\usepackage{tikz}
\usetikzlibrary{arrows.meta,positioning,shapes.geometric,calc}

% Metropolis's palette, with the decks' own colour names mapped onto it.
\definecolor{mDarkTeal}{HTML}{23373B}
\definecolor{mDarkBrown}{HTML}{604C38}
\definecolor{mLightBrown}{HTML}{EB811B}
\colorlet{crust}{mDarkBrown}
\definecolor{crumb}{HTML}{FAEEE2}
\colorlet{rye}{mDarkTeal}
\definecolor{flag}{HTML}{B03A2E}
\definecolor{own}{HTML}{2E7D4F}
\definecolor{muted}{HTML}{7A7F80}

% Top-align frames, so the \vfill in \source puts the source line at the bottom.
% A slimmer title bar and narrower margins than Metropolis's defaults leave room
% for the denser slides.
\makeatletter
\beamer@centeredfalse
\setlength{\metropolis@frametitle@padding}{1.4ex}
% A progress bar you can see: under every frame title and on section pages.
\setlength{\metropolis@progressinheadfoot@linewidth}{2pt}
\setlength{\metropolis@progressonsectionpage@linewidth}{1.5pt}
\setlength{\metropolis@titleseparator@linewidth}{1pt}
% Metropolis's title page is \paperheight tall, which overflows a top-aligned frame.
\setbeamertemplate{title page}{%
  \begin{minipage}[b][\textheight]{\textwidth}
    \vfill
    \usebeamertemplate*{title}
    \ifx\insertsubtitle\@empty\else\usebeamertemplate*{subtitle}\fi
    \usebeamertemplate*{title separator}
    \usebeamertemplate*{author}
    \usebeamertemplate*{institute}
    \vfill
  \end{minipage}}
\makeatother
\setbeamersize{text margin left=1.2em,text margin right=1.2em}

\setbeamertemplate{frame footer}{\insertshorttitle}
\setbeamertemplate{footline}{%
  \begin{beamercolorbox}[wd=\paperwidth,ht=2.6ex,dp=1.2ex,leftskip=1.2em,rightskip=1.2em]{footline}%
    \usebeamerfont{page number in head/foot}%
    \usebeamertemplate*{frame footer}\hfill\usebeamertemplate*{frame numbering}%
  \end{beamercolorbox}}
\setbeamerfont{page number in head/foot}{size=\tiny}
\setbeamercolor{footline}{fg=muted}

% The practical point of a slide, for the baker.
\newcommand{\takeaway}[1]{%
  \par\medskip
  \colorbox{crumb}{\parbox{\dimexpr\linewidth-2\fboxsep}{\small\textbf{\color{crust}In the kitchen:} #1}}}

% Small source line, placed at the bottom of a frame.
\newcommand{\source}[1]{%
  \vfill{\color{muted}\tiny Source: #1}}

% Status flags, mirroring the notes' conventions. Filled tags, so they read on the
% dark title bar as well as on the slide body.
\newcommand{\flagtag}[2]{\mbox{\setlength{\fboxsep}{1.2pt}\colorbox{#1}{\color{white}\scriptsize\bfseries\,#2\,}}}
\newcommand{\hypothesis}{\flagtag{flag}{hypothesis}}
\newcommand{\contested}{\flagtag{flag}{contested}}
\newcommand{\ownresult}{\flagtag{own}{own result}}
\newcommand{\untested}{\flagtag{flag}{not yet tested}}

% Source lists: one source per bullet, author and year in bold, optional group
% headings. Use in a [allowframebreaks] frame; long lists continue on "(cont.)" slides.
\newenvironment{sourcelist}{%
  \setbeamerfont{itemize/enumerate body}{size=\scriptsize}%
  \setbeamerfont{itemize item}{size=\scriptsize}%
  \scriptsize\begin{itemize}\setlength{\itemsep}{2pt}}{\end{itemize}}
\newcommand{\src}[2]{\item\textbf{#1}\enspace #2}
% \nobreak keeps a group heading on the same slide as its first source.
\newcommand{\srcgroup}[1]{\par\smallskip{\small\bfseries\color{crust}#1}\par\nobreak\vspace{-2pt}\nobreak}
\setbeamertemplate{frametitle continuation}[from second][(cont.)]

% Flowchart nodes for TikZ diagrams.
\tikzset{
  box/.style={draw=crust,fill=crumb,rounded corners=2pt,align=center,
    font=\scriptsize,inner sep=4pt,text width=#1},
  box/.default=2.6cm,
  arrow/.style={-{Stealth[length=5pt]},draw=crust,thick},
}

\author{Sourdough Science}
\institute{github.com/zrrbite/sourdough-science}
\date{}


% Local styles for this deck only.
\usetikzlibrary{fit,backgrounds}
\setbeamerfont{itemize/enumerate body}{size=\small}
\setbeamerfont{itemize/enumerate subbody}{size=\footnotesize}
\newcolumntype{L}{>{\raggedright\arraybackslash}X}
\newcolumntype{P}[1]{>{\raggedright\arraybackslash}p{#1}}
\tikzset{
  hot/.style={box=#1,draw=flag,fill=flag!8},
  hot/.default=2.6cm,
  cool/.style={box=#1,draw=rye!45,fill=muted!12},
  cool/.default=2.6cm,
  panel/.style={draw=muted,dashed,rounded corners=3pt,inner sep=6pt},
  plabel/.style={font=\scriptsize\bfseries,text=crust},
  elabel/.style={font=\tiny,text=muted,align=center},
}
% Own kitchen results get their own green box, so they never read as literature.
\newcommand{\ownbox}[1]{\par\smallskip{\setlength{\fboxsep}{5pt}%
  \colorbox{own!10}{\parbox{\dimexpr\linewidth-10pt}{\small\ownresult\ #1}}}\par}
\newcommand{\muted}[1]{{\color{muted}#1}}
% Tables: a group (not braces, which swallow tabularx inside a frame).
\newcommand{\tablestart}[2]{\begingroup\renewcommand{\arraystretch}{#1}#2}
\newcommand{\tableend}{\endgroup\par}

\title[The Science]{The Science: Gluten, Heat and Survival}
\subtitle{What acid, enzymes and heat do to a dough and a culture}

\begin{document}

\begin{frame}[plain]\titlepage\end{frame}

% ---------------------------------------------------------------
\begin{frame}{What this deck covers}
  Seven practical questions, each answered with its mechanism and its source:
  \begin{enumerate}\setlength{\itemsep}{0pt}
    \item Why does sourdough make dough slack and open, and when does it turn fragile?
    \item Do you need to knead and fold, or will time and the fridge do the work?
    \item Can heat buy flavour, and where does it start to cost leavening?
    \item How can a pure sourdough rise after hours at a 40°C+ setpoint?
    \item How hot is the dough really, when the oven says 45°C?
    \item How long can a culture take the heat, and how do you rescue one that didn't?
    \item Is a hot-kept culture special, and which test would settle it?
  \end{enumerate}
  \smallskip
  {\footnotesize\muted{Flags: \ownresult\ the author's own kitchen results, kept apart
  from the literature; \hypothesis, \contested, \untested\ as in the notes.
  Temperatures are oven setpoints, not dough temperatures, unless stated.}\par}
  \source{chapters 05--11 and two notes: \emph{Why the Hot Protocol Works},
    \emph{The Levain Doubling Test}}
\end{frame}

% ===============================================================
\section{Gluten}

\begin{frame}{Gluten is springs plus lubricant}
  \begin{columns}[T]
    \column{0.48\textwidth}
      {\color{crust}\textbf{Glutenin: the springs}}
      \begin{itemize}
        \item Very large polymers: subunits linked end to end by S–S (disulfide) bonds
        \item Gives \textbf{elasticity}: resists stretching and snaps back
      \end{itemize}
    \column{0.48\textwidth}
      {\color{crust}\textbf{Gliadin: the lubricant}}
      \begin{itemize}
        \item Smaller single chains, folded by internal S–S bonds only
        \item Gives \textbf{extensibility}: lets the network flow and stretch
      \end{itemize}
  \end{columns}
  \smallskip
  \begin{itemize}
    \item More glutenin: strong, snap-back dough (high-protein bread flour).
      Less: softer, more extensible (Tipo 00, pastry flour)
    \item S–S bonds keep breaking and re-forming: why dough can be reshaped, and attacked
  \end{itemize}
  \takeaway{acid, proteases and hydration changes all shift this balance. What you feel
    as stretch versus snap-back is this ratio moving.}
  \source{ch.~05; Wieser (2007); Shewry et al.\ (2002)}
\end{frame}

% ---------------------------------------------------------------
\begin{frame}{Fermentation loosens gluten four ways at once}
  \centering
  \begin{tikzpicture}[x=1cm,y=1cm]
    \node[box=1.7cm] (F) at (0,0) {Sourdough\\fermentation};
    \node[box=3.3cm] (a) at (3.1, 0.99) {pH falls to 3.5–4.5};
    \node[box=3.3cm] (b) at (3.1, 0.33) {LAB + flour proteases};
    \node[box=3.3cm] (c) at (3.1,-0.33) {Glutathione};
    \node[box=3.3cm] (d) at (3.1,-0.99) {Uneven gas production};
    \node[box=3.6cm] (a2) at (7.2, 0.99) {Chains repel: stretch ↑};
    \node[box=3.6cm] (b2) at (7.2, 0.33) {Chains cut: elasticity ↓};
    \node[box=3.6cm] (c2) at (7.2,-0.33) {S–S cross-links undone};
    \node[box=3.6cm] (d2) at (7.2,-0.99) {Patchy, uneven bubbles};
    \node[box=3.1cm,fill=white,very thick] (r) at (11.35,0)
      {\textbf{Synergistic:}\\open, irregular crumb\\+ dramatic ear};
    \foreach \n in {a,b,c,d} {
      \draw[arrow] (F.east) -- (\n.west);
      \draw[arrow] (\n) -- (\n2);
      \draw[arrow] (\n2.east) -- (r.west);
    }
  \end{tikzpicture}
  \raggedright
  \begin{itemize}
    \item Acid scales with pH: modest at pH 5.0 (10–15\% levain hybrid), dramatic at 3.8
    \item Proteases keep cutting, and flour proteases work faster at pH 3.5–4.5 than at
      neutral: modest after a 4\,h bulk, extensive after 12\,h bulk + 48\,h retard
    \item \contested\ Several studies find added acid firms dough; slackening is the
      fermentation case
  \end{itemize}
  \takeaway{levain \% sets how slack the dough gets; fermentation time decides when
    extensible tips into fragile. For pizza, that line is thin.}
  \source{ch.~05; Clarke et al.\ (2004); Thiele, Grassl \& Gänzle (2004); acidification
    note in ch.~25}
\end{frame}

% ---------------------------------------------------------------
\begin{frame}{Time builds gluten without kneading}
  \tablestart{1.0}{\footnotesize}
  \begin{tabularx}{\linewidth}{@{}P{2.0cm}LP{4.4cm}@{}}
    \toprule
    \textbf{Mechanism} & \textbf{What it does} & \textbf{Time} \\
    \midrule
    Diffusion & Thiol groups meet and form S–S bonds, as in kneading, only slower
      & 12–18\,h at ≈22°C (no-knead); 3–5× slower at 3–5°C \\
    Oxidation & Air folded in turns –SH into S–S all through the retard (what ascorbic
      acid and bromate speed up) & Builds over 3–4 days \\
    Gas pressure & Each growing bubble stretches the gluten around it: a stretch-and-fold
      from inside, everywhere at once & Continuous, even in the fridge; possibly the
      largest contributor \\
    Hydration & Water slowly unfolds the proteins, exposing buried bonding sites
      & More complete after 3–4 days \\
    \bottomrule
  \end{tabularx}
  \tableend
  {\footnotesize\muted{Kneading and folding only \emph{accelerate} these reactions, and
  the slow network is arguably more uniform.}\par}
  \takeaway{one coil fold initialises the network: it evens out hydration, adds air for
    oxidation and gives an alignment bias. For a 72–96\,h retard, one fold is enough.}
  \source{ch.~06}
\end{frame}

% ---------------------------------------------------------------
\begin{frame}{Cold slows demolition more than construction}
  \centering
  \begin{tikzpicture}[x=1cm,y=1cm]
    \node[box=2.4cm] (rc) at (0,0.5) {Construction\\≈100\% rate};
    \node[box=2.4cm] (rd) at (3.1,0.5) {Demolition\\≈100\% rate};
    \node[hot=5.5cm] (rr) at (1.55,-0.72) {After 18\,h+: demolition wins\\→ dough degrades};
    \draw[arrow] (rc.south) -- (rc.south |- rr.north);
    \draw[arrow] (rd.south) -- (rd.south |- rr.north);
    \node[box=2.4cm] (cc) at (7.6,0.5) {Construction\\20–30\% rate};
    \node[box=2.4cm] (cd) at (10.7,0.5) {Demolition\\5–10\% rate};
    \node[box=5.5cm,fill=own!12,draw=own] (cr) at (9.15,-0.72)
      {Over 3–4 days: construction wins\\→ strong network};
    \draw[arrow] (cc.south) -- (cc.south |- cr.north);
    \draw[arrow] (cd.south) -- (cd.south |- cr.north);
    \begin{scope}[on background layer]
      \node[panel,fit=(rc)(rd)(rr),label={[plabel]above:Room temperature · 22°C}] {};
      \node[panel,fit=(cc)(cd)(cr),label={[plabel]above:Cold retard · 3–5°C}] {};
    \end{scope}
  \end{tikzpicture}
  \raggedright
  \begin{itemize}
    \item Building (diffusion, oxidation, gas pressure) has low activation energies and
      keeps going in the cold; demolition (proteases, glutathione) has higher ones, so the
      cold suppresses it disproportionately. The rates are rough estimates
  \end{itemize}
  \takeaway{mix, fold once, ball and retard for 48–72\,h. Hands-on time drops from
    3–4\,h to about 30\,min, and the cold builds the network as it develops flavour.}
  \source{ch.~06}
\end{frame}

% ===============================================================
\section{Heat as a flavour tool}

\begin{frame}{Heat-stressed yeast make more flavour compounds, not fewer}
  \begin{itemize}
    \item Usual ranges (levain 25–28°C, bulk 22–25°C) target schedule and leavening, not flavour
    \item Near their stress threshold (roughly 35–40°C for most strains) yeast mount a
      \textbf{heat shock response}: chaperone proteins (HSPs), trehalose, more
      saturated membranes
    \item Fusel alcohols rise (isoamyl alcohol, isobutanol, 2-phenylethanol): the
      precursors of fruity, floral esters. The ATF1/ATF2 ester enzymes also run faster
    \item LAB thrive where yeast are stressed (\textit{L.~delbrueckii} grows best at
      42–45°C; many lactobacilli fully active at 35–40°C): more acid, proteolysis, LAB aromas
    \item Net: more flavour per hour, in a shorter window than 8–12\,h at 25°C
  \end{itemize}
  \takeaway{a warm levain is a flavour lever, not just a speed-up. The next slide is
    the catch.}
  \source{ch.~07; Verstrepen et al.\ (2003); Saerens et al.\ (2010); Kucharczyk \&
    Tuszyński (2018); Landaud et al.\ (2001)}
\end{frame}

% ---------------------------------------------------------------
\begin{frame}{Hotter does not reliably mean more esters \contested}
  \begin{columns}[T]
    \column{0.48\textwidth}
      {\color{crust}\textbf{More with heat: brewing}}
      \begin{itemize}\footnotesize
        \item Kucharczyk \& Tuszyński (2018): esters and higher alcohols rose
          significantly with temperature, at industrial scale
        \item Landaud et al.\ (2001): temperature drives fusel alcohols independently
          of pressure
      \end{itemize}
    \column{0.48\textwidth}
      {\color{crust}\textbf{Fewer with heat: sourdough yeast}}
      \begin{itemize}\footnotesize
        \item Ferrández-Gómez et al.\ (2024): sourdough strains in beer wort at 37°C gave
          a \emph{drastic reduction} in volatiles, esters especially
      \end{itemize}
  \end{columns}
  \medskip
  \begin{itemize}
    \item Likely reason: heat speeds both ester synthesis and volatile loss. Which wins
      depends on strain, headspace, surface-to-volume and duration
    \item Brief, closed, dough-bound: likely a net gain. Long, open, runny: likely a loss
  \end{itemize}
  \takeaway{don't assume hotter means fruitier: taste it against a room-temperature
    control. That side-by-side is still to be done. \untested}
  \source{ch.~07, caveat on conflicting evidence}
\end{frame}

% ---------------------------------------------------------------
\begin{frame}{The thermal risk gradient is a slope, not a cliff}
  \tablestart{1.0}{\footnotesize}
  \begin{tabularx}{\linewidth}{@{}P{1.3cm}LP{3.7cm}P{3.4cm}@{}}
    \toprule
    & \textbf{Yeast} & \textbf{LAB} & \textbf{For the baker} \\
    \midrule
    25–30°C & Comfort zone; no stress response & Active & Most recipes \\
    30–35°C & Near optimum; mild stress at the top & Accelerate in step
      & Faster; most cope \\
    \textbf{35–40°C} & Full heat shock response; growth slows, metabolites rise
      & Thriving; may dominate & \textbf{The flavour sweet spot} \\
    40–45°C & Most strains lose viability toward the top; a gradual decline
      & Thermophiles may stay fully active & Flavour; gas may fail \\
    45°C+ & Thermal kill for most yeast & Mesophiles decline; thermophiles may persist
      & Not useful for dough \\
    \bottomrule
  \end{tabularx}
  \tableend
  {\footnotesize\muted{Strain- and culture-dependent, from lab studies: a warm-kept or
  heat-selected culture may exceed these.}\par}
  \takeaway{if a 40–45°C build leavens sluggishly, the yeast have fallen below functional
    levels. The LAB still make flavour, so it tastes right but won't rise.}
  \source{ch.~07, from the literature on the most studied sourdough species}
\end{frame}

% ---------------------------------------------------------------
\begin{frame}{Hot sprint, cold coast: make flavour hot, keep it cold}
  \centering
  \begin{tikzpicture}
    \node[hot=4.6cm] (hs) {\textbf{Hot sprint · 35–40°C}\\[2pt]
      heat shock response\\fusel alcohols ↑ · LAB flavour ↑\\fast acid production};
    \node[cool=4.6cm,right=2.6cm of hs] (cc) {\textbf{Cold coast · 3–5°C, 48–72\,h}\\[2pt]
      volatiles preserved\\gluten builds · sugars depleted\\Maillard precursors ↑};
    \draw[arrow] (hs) -- node[above,font=\scriptsize]{precursor pool}
      node[below,font=\scriptsize]{locked in} (cc);
  \end{tikzpicture}
  \raggedright
  \begin{itemize}
    \item The catch: more LAB activity means more proteases, so a hot front end loads
      gluten damage onto the dough before the fridge
    \item Limit it: keep the \emph{dough} hot 30–60\,min at most; use 10\% levain rather
      than 15–20\%; use a high-W flour (Caputo Oro W 350–380+). W 280–310 may not survive
  \end{itemize}
  \takeaway{build the levain at 35–38°C: it peaks in 4–6\,h rather than 8–12, and passes
    peak faster, so catch it right at peak. Mix the dough at ≈22°C, fold once, retard.}
  \source{ch.~07}
\end{frame}

% ===============================================================
\section{The hot culture and Kveik}

\begin{frame}{A pure sourdough rose after ≈3\,h at a 40°C+ setpoint}
  \centering
  \begin{tikzpicture}[node distance=0.55cm]
    \node[hot=2.7cm] (s1) {\textbf{1 · Levain}\\40°C+, ≈1\,h};
    \node[hot=2.7cm,right=of s1] (s2) {\textbf{2 · Bulk}\\40°C+, ≈2\,h};
    \node[cool=2.7cm,right=of s2] (s3) {\textbf{3 · Cold retard}\\3–5°C, 10–12\,h};
    \node[box=2.7cm,fill=own!12,draw=own,right=of s3] (s4)
      {\textbf{4 · Bake}\\good rise, no failure};
    \draw[arrow] (s1) -- (s2); \draw[arrow] (s2) -- (s3); \draw[arrow] (s3) -- (s4);
  \end{tikzpicture}
  \raggedright
  \ownbox{Repeated with success: well-leavened crumb, no structural failure, and no
    commercial yeast. All the gas came from the culture's wild yeast.}
  \begin{itemize}
    \item The literature predicts inhibition or death here: \textit{K.~humilis}, a
      textbook sourdough yeast, is inhibited at 35°C (Brandt et al.\ 2004)
    \item 40°C+ is the oven setting. The dough itself was never probed
  \end{itemize}
  \takeaway{test your own culture rather than trusting species data. If a 40°C+ build
    gives a well-risen loaf, it is heat-tolerant enough for that protocol.}
  \source{ch.~08 (own kitchen results); Brandt et al.\ (2004)}
\end{frame}

% ---------------------------------------------------------------
\begin{frame}{Two explanations, and they fail in different ways}
  \tablestart{1.05}{\footnotesize}
  \begin{tabularx}{\linewidth}{@{}P{2.6cm}LL@{}}
    \toprule
    & \textbf{1 · Genuine thermotolerance} \hypothesis
    & \textbf{2 · Temporal escape} \hypothesis \\
    \midrule
    The idea & The culture holds heat-adapted yeast, active and dividing at 40°C
      & Ordinary yeast, genuinely stressed; ≈3\,h is short enough that enough survive \\
    Gas comes from & Yeast still growing at temperature
      & Non-dividing cells still fermenting sugar, plus CO₂ dissolved in the cold \\
    Why plausible & \textit{Pichia kudriavzevii} survives 40°C in sourdough (Meroth et
      al.\ 2003); warm builds could select, as with Kveik
      & Heat kill is time-dependent: at 40°C, gradual rather than instant \\
    Hot bulk stretched to 4–8\,h & Still makes gas (gluten may fail first)
      & Leavening collapses abruptly, probably between 4 and 8\,h \\
    \bottomrule
  \end{tabularx}
  \tableend
  \smallskip
  \begin{itemize}
    \item \untested\ \textbf{Open.} Neither the extended-bulk test nor the levain doubling
      test has been run. A note bets on explanation 2 (last section)
  \end{itemize}
  \source{ch.~08; Meroth et al.\ (2003); Brandt et al.\ (2004)}
\end{frame}

% ---------------------------------------------------------------
\begin{frame}{Warm practice selects microbes, over centuries or days}
  \tablestart{1.0}{\footnotesize}
  \begin{tabularx}{\linewidth}{@{}P{3.4cm}L@{}}
    \toprule
    \textbf{Study} & \textbf{What it found} \\
    \midrule
    \multicolumn{2}{@{}l}{\scriptsize\color{muted}KVEIK: farmhouse yeast pitched warm
      (often 28–40°C+) and reused from good batches} \\
    Preiss et al.\ (2018) & Kveik from nine Norwegian sources: a genetically distinct
      group of domesticated \textit{S.~cerevisiae}, with marked ethanol and heat tolerance \\
    Foster et al.\ (2022) & More trehalose than standard ale yeast: a direct
      heat-tolerance mechanism \\
    \midrule
    \multicolumn{2}{@{}l}{\scriptsize\color{muted}SOURDOUGH: same backslopping,
      different temperatures} \\
    Van der Meulen et al.\ (2011) & Wheat sourdough: 23°C selected
      \textit{Leuconostoc citreum}; 30 and 37°C selected \textit{L.~fermentum}, a
      complete shift within 10 backslopping cycles \\
    Urien et al.\ (2014) & Rye sourdough, 56 daily cycles: at 20°C
      \textit{Kazachstania unispora} alone; at 30°C \textit{S.~cerevisiae},
      \textit{P.~kudriavzevii} and \textit{C.~glabrata} \\
    \bottomrule
  \end{tabularx}
  \tableend
  \takeaway{your feeding temperature is a selection pressure: on its own it can replace
    a sourdough's dominant species in about ten feeds.}
  \source{ch.~09; Preiss et al.\ (2018); Foster et al.\ (2022); Van der Meulen et al.\
    (2011); Urien et al.\ (2014)}
\end{frame}

% ---------------------------------------------------------------
\begin{frame}{At home, heat selects species, not new genes}
  \centering
  \begin{tikzpicture}[x=1cm,y=1cm]
    \node[hot=1.5cm] (A) at (0,0) {Repeated hot builds, 40°C+};
    \node[box=2.3cm] (B) at (2.63,0) {Heat-tolerant cells keep dividing; mesophiles
      stall};
    \node[box=2.3cm] (C) at (5.61,0) {Competitive exclusion: species shift in
      ≈10 cycles};
    \node[box=3.6cm] (E) at (9.29,0.72) {\textbf{Yeast} → \textit{P.~kudriavzevii}
      (to 45°C), wild \textit{S.~cerevisiae}, \textit{W.~anomalus}};
    \node[box=3.6cm] (G) at (9.29,-0.72) {\textbf{LAB} → from
      \textit{F.~sanfranciscensis} (to ≈40°C) to \textit{L.~delbrueckii},
      \textit{helveticus}, \textit{fermentum}};
    \node[box=2.2cm,fill=white,very thick] (I) at (12.85,0.85) {Species-level
      selection,\\months to years};
    \node[cool=2.2cm] (K) at (12.85,-0.95) {\textbf{Kveik}: centuries,\\plus genomic
      change};
    \draw[arrow] (A) -- (B); \draw[arrow] (B) -- (C);
    \draw[arrow] (C.east) -- (E.west); \draw[arrow] (C.east) -- (G.west);
    \draw[arrow] (E.east) -- (I.west); \draw[arrow] (G.east) -- (I.west);
    \draw[arrow,dashed] (I) -- (K);
  \end{tikzpicture}
  \raggedright
  \begin{itemize}
    \item The LAB shift favours lactic over acetic: may suit pizza, as acetic tightens dough
    \item \hypothesis\ That this has happened to the author's culture is ch.~08's open
      explanation 1
  \end{itemize}
  \takeaway{selection only bites if you backslop from hot builds: 10–20 cycles reshape
    the culture. A mother kept at 20–25°C keeps its full consortium, so keep one as backup.}
  \source{ch.~08–09; Van der Meulen et al.\ (2011); Meroth et al.\ (2003); Limtong et al.\
    (2018); Brandt et al.\ (2004)}
\end{frame}

% ===============================================================
\section{Thermal lag}

\begin{frame}{The oven setpoint is not the dough temperature}
  Dough's thermal diffusivity is like water's (≈1.4\,×\,10\textsuperscript{−7}\,m²/s);
  plastic insulates further.\par\smallskip
  \tablestart{1.0}{\footnotesize}
  \begin{tabularx}{\linewidth}{@{}P{3.6cm}L@{}}
    \toprule
    \textbf{Vessel at a 45°C setpoint} & \textbf{How long the centre lags} \\
    \midrule
    Small jar, 75\,g & Warms through in a small fraction of an hour \\
    Levain jar, 200–300\,g & May take 30–45\,min to reach setpoint \\
    Dough tub, 1\,kg+ & An hour or more \\
    \bottomrule
  \end{tabularx}
  \tableend
  \begin{itemize}
    \item ``1\,h at 40°C+'', 200–300\,g jar: core 25–35°C for the first half; average
      probably 30–38°C
    \item \untested\ At 40–50°C settings the core is likely 5–15°C below setpoint most of
      the time: an estimate, as no dough has been probed yet
  \end{itemize}
  \takeaway{probe the centre of the dough and log it through a build: time to setpoint,
    time-weighted average, and the surface-to-core gap.}
  \source{ch.~10; Sadot et al.\ (2017); 75\,g jar: \emph{The Levain Doubling Test};
    oven air measured ≈45°C on the usual setting (2026-09-13)}
\end{frame}

% ---------------------------------------------------------------
\begin{frame}{One dough, two fermentations: hot shell, cool core}
  \centering
  \begin{tikzpicture}[x=1cm,y=1cm]
    \node[box=2.1cm] (S) at (0,0) {Setpoint 45°C;\\dough from 22–24°C;\\core lags};
    \node[hot=3.0cm] (O) at (3.28,0.62) {\textbf{Outer layer · 40–45°C}\\LAB at peak;
      yeast stressed};
    \node[cool=3.0cm] (C) at (3.28,-0.62) {\textbf{Core · 30–35°C}\\wild yeast dividing};
    \node[box=3.4cm] (OF) at (7.21,0.62) {Flavour reactor:\\acids, proteases, diacetyl};
    \node[box=3.4cm] (CF) at (7.21,-0.62) {Leavening zone:\\CO₂, fusel alcohols};
    \node[box=3.0cm,fill=white,very thick] (Z) at (11.14,0) {Flavour \emph{and}\\
      leavening together};
    \draw[arrow] (S.east) -- (O.west); \draw[arrow] (S.east) -- (C.west);
    \draw[arrow] (O) -- (OF); \draw[arrow] (C) -- (CF);
    \draw[arrow] (OF.east) -- (Z.west); \draw[arrow] (CF.east) -- (Z.west);
  \end{tikzpicture}
  \raggedright
  \begin{itemize}
    \item A uniform 40°C would cost leavening; a uniform 30°C would cost flavour
    \item The yeast may never feel what the dial implies, so heat selection is softer than
      the setpoint suggests (the shell is still above 40°C)
  \end{itemize}
  \takeaway{if the unstirred hot build works, don't stir: it collapses the gradient, and
    its oxygen grows more yeast but fewer fusels and esters. To test, stir once at halfway.}
  \source{ch.~10}
\end{frame}

% ===============================================================
\section{Thermal death and recovery}

\begin{frame}{Thermal death is a rate, not a threshold}
  \ownbox{Brief excursions to 40°C+ (≈1\,h levain, ≈2\,h bulk, each then retarded cold)
    did no apparent harm over several days. One prolonged session at the same temperature
    killed, or very nearly killed, the culture.}
  \begin{itemize}
    \item Heat kills at a rate: numbers fall exponentially with time at temperature
    \item \textbf{D-value}: time for a 90\% (one-log) kill. \textbf{z-value}: the rise
      that cuts D tenfold
    \item So 10\,min at 45°C and 2\,h at 42°C may kill about equally
    \item Lab D-values for \textit{S.~cerevisiae}: 22.3 down to 0.8\,min across 45–55°C,
      in buffer or media. In dough they are lower bounds on survival: use qualitatively
  \end{itemize}
  \takeaway{a temperature on its own says nothing about survival. Always give the time
    at temperature too.}
  \source{ch.~11; Guerrero, López-Malo \& Alzamora (1999/2001)}
\end{frame}

% ---------------------------------------------------------------
\begin{frame}{The heat-shock defence is a budget, and the fridge refills it}
  \begin{itemize}
    \item It works (a mild 37°C pre-shock induces tolerance to otherwise lethal heat;
      Morano et al.\ 2012) but burns ATP, uses up trehalose, and repairs damage without
      preventing it
  \end{itemize}
  \medskip
  \centering
  \begin{tikzpicture}[x=1cm,y=1cm]
    \node[plabel,anchor=east,align=right] at (-0.15,0.6) {Brief, repeated\\≈1–2\,h};
    \node[hot=2.0cm] (a) at (1.2,0.6) {Hot session\\40°C+};
    \node[box=2.0cm] (b) at (3.75,0.6) {Heat shock response};
    \node[cool=2.0cm] (c) at (6.3,0.6) {Cold retard:\\recovery window};
    \node[box=2.0cm] (d) at (8.85,0.6) {Rebuild proteins, trehalose, ATP};
    \node[box=1.9cm,fill=own!12,draw=own] (e) at (11.4,0.6) {Dips, recovers:
      sustainable};
    \draw[arrow] (a) -- (b); \draw[arrow] (b) -- (c); \draw[arrow] (c) -- (d);
    \draw[arrow] (d.north) -- ++(0,0.25) -| (a.north);
    \draw[arrow] (d) -- (e);
    \node[plabel,anchor=east,align=right] at (-0.15,-0.95) {One long\\session};
    \node[hot=2.0cm] (f) at (1.2,-0.95) {Continuous stress, no recovery};
    \node[box=2.0cm] (g) at (3.75,-0.95) {ATP drained, trehalose spent};
    \node[box=2.0cm] (h) at (6.3,-0.95) {Repair falls behind damage};
    \node[hot=2.0cm] (i) at (8.85,-0.95) {Stressed → dead, irreversible};
    \node[box=1.9cm] (j) at (11.4,-0.95) {Defence runs to depletion};
    \draw[arrow] (f) -- (g); \draw[arrow] (g) -- (h); \draw[arrow] (h) -- (i);
    \draw[arrow] (i) -- (j);
  \end{tikzpicture}
  \raggedright
  \takeaway{the cold retard after each hot session is part of the method, not a pause:
    it is where the culture recovers.}
  \source{ch.~11; Morano et al.\ (2012)}
\end{frame}

% ---------------------------------------------------------------
\begin{frame}{Two free protections: a slow ramp and a stiff levain}
  \begin{columns}[T]
    \column{0.48\textwidth}
      {\color{crust}\textbf{Ramp, not shock}}
      \begin{itemize}\footnotesize\setlength{\itemsep}{1pt}
        \item \textit{S.~cerevisiae} survives a later hour at 50°C considerably better
          after a slow, linear warm-up than after a sudden shock
        \item Gained mainly between 40 and 50°C on the way up, by membrane remodelling,
          not Hsp104 or trehalose
        \item So the container's lag protects: pre-warming it, or dropping a small levain
          into a hot water bath, is harsher than it looks
      \end{itemize}
    \column{0.48\textwidth}
      {\color{crust}\textbf{Lower water activity}}
      \begin{itemize}\footnotesize\setlength{\itemsep}{1pt}
        \item Heat resistance rises as water activity falls: a stiff levain does from
          outside what trehalose does inside
        \item 100\% vs 50\% hydration is perhaps a\textsubscript{w} 0.99 vs 0.95: plausibly
          a factor of a few in D-value. \untested\ No published value at 40–45°C
        \item 100\% to 85\% buys very little; a true 50–60\% stiff levain buys more
      \end{itemize}
  \end{columns}
  \takeaway{start hot builds from cool, in their container, and stiffen the levain
    \emph{before} raising the heat. A margin, not a licence: neither extends a safe
    session.}
  \source{ch.~11; Guyot, Ferret \& Gervais; Morano et al.\ (2012)}
\end{frame}

% ---------------------------------------------------------------
\begin{frame}{Long sessions are harsher: acid, ethanol and hunger compound}
  \begin{columns}[T]
    \column{0.44\textwidth}
      \begin{itemize}\footnotesize
        \item \textbf{Falling pH} speeds heat kill; pH control and heat defence share
          one ATP budget
        \item \textbf{Ethanol} speeds heat kill exponentially; both hit the membrane, so
          they multiply (Leão \& van Uden 1982; Piper 1995)
        \item \textbf{Sugar runs out}: no fuel for the defence, no carbon to rebuild
          trehalose
      \end{itemize}
      \smallskip
      {\footnotesize\itshape ``Yeast that comfortably survives 45°C at pH 5.5 with ample
      maltose may die at 45°C at pH 3.8 with none.''\par}
    \column{0.54\textwidth}
      \tablestart{1.2}{\scriptsize}
      \begin{tabularx}{\linewidth}{@{}P{2.7cm}L@{}}
        \toprule
        \textbf{Regime} & \textbf{Assessment} \\
        \midrule
        Brief (1–3\,h) at 40–45°C, recovery between & Sustainable, mildly selective:
          the documented working protocol \\
        Extended (6\,h+) at 40–45°C, no recovery & High risk: the defence exhausts,
          stressors pile up \\
        Any duration at 50°C+ & Thermal kill for most sourdough yeast \\
        Hot session on an acidic, sugar-depleted culture & Substantially more lethal
          than on a fresh feed \\
        \bottomrule
      \end{tabularx}
      \tableend
  \end{columns}
  \takeaway{feed immediately \emph{before} a hot session, not after, and build stiff.
    Neither extends the time a culture can take; both raise its odds.}
  \source{ch.~11; Leão \& van Uden (1982); Piper (1995)}
\end{frame}

% ---------------------------------------------------------------
\begin{frame}{Sour but flat means the yeast died first}
  \begin{columns}[T]
    \column{0.50\textwidth}
      \tablestart{1.2}{\scriptsize}
      \begin{tabularx}{\linewidth}{@{}P{2.6cm}L@{}}
        \toprule
        \textbf{You see} & \textbf{It means} \\
        \midrule
        Sharp acid, no gas, no rise & \textbf{Yeast lost, LAB surviving.} The most
          common outcome \\
        No acid, no gas, flat or ``off'' smell & Both badly damaged; poorer prognosis \\
        Gas back within 1–2 feeds, acid intact & Reduced but viable; recovering \\
        Sweet or solvent smell, no gas, slimy & Possible spoilage: discard \\
        \bottomrule
      \end{tabularx}
      \tableend
      \smallskip
      {\footnotesize LAB are considerably more heat-tolerant, so a culture can smell
      healthy and sour with no leavening left.\par}
    \column{0.47\textwidth}
      {\color{crust}\textbf{Recovery}: a 4–5 log kill leaves survivors}\par
      \vspace{-0.6em}
      \begin{enumerate}\footnotesize\setlength{\itemsep}{1pt}
        \item Feed at 22–25°C: warmth was the stressor
        \item Feed ≥1:10:10 to dilute acid and ethanol
        \item Add wholemeal or high-extraction flour
        \item Feed daily for 3–7 days before judging
        \item Pink, orange, fuzzy or putrid: discard
      \end{enumerate}
  \end{columns}
  \takeaway{keep the recovered culture as a second culture, not a replacement, and
    compare it with the original or a room-temperature backup.}
  \source{ch.~11}
\end{frame}

% ===============================================================
\section{Why the hot protocol works}

\begin{frame}{At 28°C nothing is stressed: the usual advice is about schedule}
  \begin{itemize}\setlength{\itemsep}{0.3em}
    \item \textit{S.~cerevisiae} grows best around 30–33°C; most sourdough LAB at
      30–40°C. A levain at 28°C is near an \textbf{optimum}, not a limit
    \item Why the usual advice says 24–26°C:
      \begin{itemize}
        \item \textbf{Schedule control}: cooler widens the timing window and forgives a
          late baker. The largest single reason
        \item \textbf{A flavour preference stated as a rule}: warmer means more acid per
          unit of gas, and more lactic than acetic. A tradition chasing acetic sharpness
          runs cool
      \end{itemize}
    \item It is not even universal: Italian \emph{lievito madre} is classically kept
      around 28°C, and panettone practice runs at 26–30°C
  \end{itemize}
  \takeaway{at 28°C you trade a narrower timing window and a rounder, more lactic sour.
    Not viability.}
  \source{\emph{Why the Hot Protocol Works}, part 1; ch.~07}
\end{frame}

% ---------------------------------------------------------------
\begin{frame}{Six ordinary effects compound: no exotic culture needed \hypothesis}
  \centering
  \begin{tikzpicture}[x=1cm,y=1cm,eff/.style={box=3.0cm,inner sep=3pt,minimum height=1.15cm}]
    \node[eff] (e1) at (3.1,0) {\textbf{1 · Thermal lag}\\core below setpoint;
      mean probably low-to-mid 30s};
    \node[eff] (e2) at (6.5,0) {\textbf{2 · Short, then cold}\\≈1\,h levain,
      ≈2\,h bulk, then fridge};
    \node[eff] (e3) at (9.9,0) {\textbf{3 · Fed just before}\\most sugar, least
      acid and ethanol};
    \node[eff,anchor=north] (e4) at ([yshift=-4pt]e1.south) {\textbf{4 · Slow ramp}\\not
      shock: membranes remodel on the climb};
    \node[eff] (e5) at (e2 |- e4) {\textbf{5 · Low water activity}\\a stiff levain
      protects more};
    \node[eff] (e6) at (e3 |- e4) {\textbf{6 · Cool-core refuge}\\shell makes
      flavour, core makes gas};
    \begin{scope}[on background layer]
      \node[panel,inner sep=3pt,fit=(e1)(e2)(e3)(e4)(e5)(e6)] (stack) {};
    \end{scope}
    \node[hot=1.5cm,anchor=east] (SP) at ([xshift=-0.4cm]stack.west)
      {Setpoint 45°C: what the dial says};
    \node[box=1.8cm,fill=own!12,draw=own,anchor=west] (OUT) at ([xshift=0.4cm]stack.east)
      {Well-leavened loaf, without an unusual culture};
    \draw[arrow] (SP) -- (stack.west);
    \draw[arrow] (stack.east) -- (OUT);
    \node[elabel,anchor=north] at ([yshift=-3pt]OUT.south) {multiplying,\\not adding};
  \end{tikzpicture}
  \raggedright
  \begin{itemize}
    \item \textbf{The position}: temporal escape plus five other effects does nearly all
      the work. ``It works'' is weak evidence: a 90\% kill still leaves enough to leaven
  \end{itemize}
  {\footnotesize\color{flag}\textbf{A stated bet, not a result.} Heat-tolerant selection
  is at most a small extra term; the test that would falsify the bet has not been run.\par}
  \source{\emph{Why the Hot Protocol Works}, parts 2–3; ch.~08, 10, 11}
\end{frame}

% ---------------------------------------------------------------
\begin{frame}{A doubling levain shows gas, not cell division \untested}
  \begin{columns}[T]
    \column{0.53\textwidth}
      {\footnotesize One feed split into two equal jars: one hot, one at room
      temperature.\par}
      \smallskip
      \tablestart{1.2}{\scriptsize}
      \begin{tabularx}{\linewidth}{@{}P{2.5cm}L@{}}
        \toprule
        \textbf{Hot jar does} & \textbf{It supports} \\
        \midrule
        Strong, prompt rise, like the control & Yeast active at temperature:
          \textbf{consistent with either explanation} \\
        Acid and aroma, little or no rise & LAB active, yeast not contributing gas:
          temporal escape \\
        Rise markedly slower than the control & Yeast impaired but working: temporal
          escape \\
        \bottomrule
      \end{tabularx}
      \tableend
      \smallskip
      {\footnotesize Stressed, non-dividing yeast still make CO₂. A negative result is
      informative; a positive one is \textbf{not decisive}.\par}
    \column{0.44\textwidth}
      {\color{crust}\textbf{At the bench}}
      \begin{itemize}\footnotesize
        \item Same mother, same feed, \textbf{equal mass}: mass drives the thermal lag
        \item Identical tall, narrow jars; rubber band at the start line
        \item Control at room temperature, \emph{measured} and written down
        \item Height every 15\,min (photo with a ruler); oven air each time; core probe if
          you have one
        \item Feed the mother the day before; run until the hot jar peaks and recedes
      \end{itemize}
  \end{columns}
  \source{\emph{The Levain Doubling Test}; ch.~08}
\end{frame}

% ---------------------------------------------------------------
\begin{frame}{Serial hot passage is the test that would settle it \untested}
  \begin{columns}[T]
    \column{0.56\textwidth}
      \begin{itemize}\footnotesize
        \item At a 1:1:1 feed the old starter is a third of the new mass
        \item If the yeast are not dividing, cells per gram fall to
          (1/3)\textsuperscript{\emph{n}}: a twenty-seventh by build 3, under a
          two-hundredth by build 5
        \item Five hot builds in a row, each backslopped \textbf{from the hot jar}, not
          the fridge or the mother; same ratio, time, setpoint
        \item If build 5 still rises like build 1, the yeast are dividing at
          temperature: explanation 1, and the bet is wrong
      \end{itemize}
    \column{0.41\textwidth}
      \centering
      \begin{tikzpicture}[x=0.75cm,y=2.3cm]
        \node[font=\scriptsize,align=center,text=rye] at (2.5,1.32)
          {Cells per gram if the yeast\\are \emph{not} dividing (1:1:1)};
        \draw[muted] (-0.5,0) -- (5.5,0);
        \foreach \n/\v/\l in {0/1/1, 1/0.3333/{1/3}, 2/0.1111/{1/9}, 3/0.0370/{1/27},
                              4/0.0123/{1/81}, 5/0.0041/{<\,1/200}} {
          \fill[crust!80] (\n-0.3,0) rectangle (\n+0.3,\v);
          \node[font=\tiny,above] at (\n,\v) {\l};
          \node[font=\tiny,below,text=muted] at (\n,0) {\n};
        }
        \node[font=\tiny,text=muted] at (2.5,-0.2) {build \emph{n}};
      \end{tikzpicture}
  \end{columns}
  \takeaway{split off a test jar first. Never run this on the only culture in the house.}
  \source{\emph{The Levain Doubling Test}; ch.~08}
\end{frame}

% ===============================================================
\begin{frame}{Takeaways}
  \begin{itemize}\setlength{\itemsep}{0.35em}
    \item Acid, proteases and glutathione slacken gluten: levain \% sets how much,
      fermentation time decides when extensible turns fragile
    \item Time and cold build gluten: one fold and a 48–72\,h retard replace most kneading
    \item Heat can buy flavour (35–40°C), but more esters is \emph{not} guaranteed: taste
      against a control, and keep the dough's own hot phase short
    \item The dial is not the dough: a 200–300\,g jar's centre may lag 30–45\,min, a 1\,kg
      tub an hour or more
    \item Survival is time at temperature: brief hot sessions with a cold recovery are
      sustainable, long ones high risk
    \item Feed right before the heat, start from cool, build stiff, keep a
      room-temperature backup
    \item Sour but flat means the yeast died: recover at 22–25°C, 1:10:10, for 3–7 days
  \end{itemize}
\end{frame}

% ---------------------------------------------------------------
\begin{frame}{Open questions}
  \tablestart{1.0}{\footnotesize}
  \begin{tabularx}{\linewidth}{@{}P{4.6cm}LP{2.3cm}@{}}
    \toprule
    \textbf{Question} & \textbf{Test} & \textbf{Status} \\
    \midrule
    Heat-adapted, or temporal escape? (ch.~08; a note bets on escape)
      & Serial hot passage (decisive); levain doubling; extend the hot bulk until
      leavening fails & \untested \\
    How hot is the dough really? (ch.~10) & Log the jar core in a hot build; 75\,g jar
      vs 1\,kg tub & \untested \\
    Does heat raise or lower esters? (ch.~07) & Taste side by side with a
      room-temperature control & \contested \\
    How much does a stiff levain protect? (ch.~11) & Stiff vs liquid, same setpoint and
      time; then feed and time to peak & \untested \\
    Is the recovered culture better? (ch.~11) & Against a control: heat tolerance,
      flavour, timing at 25°C & \untested \\
    Stir the hot build or not? (ch.~10) & One stir at halfway vs none: rise and flavour
      & \leavevmode\muted{suggested} \\
    \bottomrule
  \end{tabularx}
  \tableend
  \source{ch.~07, 08, 10, 11; notes; README status list. The \texttt{chamber/} logger
    meant to run these has its first stage built; nothing is instrumented yet}
\end{frame}

% ---------------------------------------------------------------
\begin{frame}[allowframebreaks]{Sources}
  \srcgroup{Gluten and acidification}
  \begin{sourcelist}
    \src{Wieser H (2007)}{Chemistry of gluten proteins. \emph{Food Microbiology} 24(2):115--119.}
    \src{Shewry PR, Halford NG, Belton PS, Tatham AS (2002)}{The structure and properties of gluten. \emph{Phil.\ Trans.\ R.\ Soc.\ B} 357:133--142.}
    \src{Clarke CI et al.\ (2004)}{Wheat sourdough fermentation: effects of time and acidification on rheological properties. \emph{Cereal Chemistry} 81(3):409--417.}
    \src{Thiele C, Grassl S, Gänzle MG (2004)}{Gluten hydrolysis and depolymerization during sourdough fermentation. \emph{J.\ Agric.\ Food Chem.} 52(5):1307--1314.}
  \end{sourcelist}
  \srcgroup{Flavour and temperature}
  \begin{sourcelist}
    \src{Verstrepen KJ et al.\ (2003)}{Expression levels of the yeast alcohol acetyltransferase genes. \emph{Appl.\ Environ.\ Microbiol.} 69(9):5228--5237.}
    \src{Saerens SMG et al.\ (2010)}{Production and biological function of volatile esters in \textit{S.~cerevisiae}. \emph{Microbial Biotechnology} 3(2):165--177.}
    \src{Landaud S, Latrille E, Corrieu G (2001)}{\emph{J.\ Inst.\ Brewing} 107(2):107--117.}
    \src{Kucharczyk K, Tuszyński T (2018)}{\emph{J.\ Inst.\ Brewing} 124(3):230--235.}
    \src{Ferrández-Gómez B et al.\ (2024)}{Sourdough yeast strains \ldots\ in beer wort. \emph{Foods} 13.}
  \end{sourcelist}
  \srcgroup{Heat transfer}
  \begin{sourcelist}
    \src{Sadot M et al.\ (2017)}{Numerical modelling of convective heat transfer during bread baking. \emph{J.\ Food Engineering}.}
  \end{sourcelist}
  \srcgroup{Selection, kveik and sourdough ecology}
  \begin{sourcelist}
    \src{Preiss R et al.\ (2018)}{Traditional Norwegian kveik are a genetically distinct group of domesticated \textit{S.~cerevisiae}. \emph{Front.\ Microbiol.} 9:2137.}
    \src{Foster B et al.\ (2022)}{Kveik brewing yeasts \ldots\ enhanced trehalose accumulation. \emph{Front.\ Microbiol.} 13:747546.}
    \src{Van der Meulen R et al.\ (2011)}{\emph{Appl.\ Environ.\ Microbiol.} 77(8):2716--2726.}
    \src{Urien C et al.\ (2014)}{\emph{PLOS ONE} 9(4):e95449.}
    \src{Meroth CB, Hammes WP, Hertel C (2003)}{\emph{Appl.\ Environ.\ Microbiol.} 69(12):7453--7461.}
    \src{Brandt MJ, Hammes WP, Gänzle MG (2004)}{\emph{Eur.\ Food Res.\ Technol.} 218:333--338.}
    \src{Limtong S, Sringiew C, Yongmanitchai W (2018)}{\emph{Agriculture and Natural Resources} 52.}
  \end{sourcelist}
  \srcgroup{Thermal death and stress}
  \begin{sourcelist}
    \src{Morano KA, Grant CM, Moye-Rowley WS (2012)}{\emph{Genetics} 190(4):1157--1195.}
    \src{Guyot S, Ferret E, Gervais P}{Thermal death rate of \textit{S.~cerevisiae}: dependence on the kinetics of heat perturbation.}
    \src{Guerrero S, López-Malo A, Alzamora SM (1999/2001)}{\textit{S.~cerevisiae} thermal inactivation kinetics combined with ultrasound.}
    \src{Leão C, van Uden N (1982)}{\emph{Biotechnol.\ Bioeng.} 24:1581--1590.}
    \src{Piper PW (1995)}{\emph{FEMS Microbiol.\ Lett.} 134(2--3):121--127.}
  \end{sourcelist}
  \medskip
  {\scriptsize\muted{Full list: \texttt{docs/25-references.md}}\par}
\end{frame}

\end{document}
